%!TEX program = xelatex
% 口語報告提示稿（中英對照）— 對應 presentation.pdf 之 20 頁，約 40 分鐘
\documentclass[10.5pt, a4paper]{article}

\usepackage{fontspec}
\usepackage{xeCJK}
\usepackage{geometry}
\usepackage{setspace}
\usepackage{titlesec}
\usepackage{amsmath}
\usepackage{amssymb}
\usepackage{xcolor}
\usepackage{fancyhdr}
\usepackage{tcolorbox}

\geometry{top=1.6cm, bottom=1.6cm, left=2.0cm, right=2.0cm}
\setlength{\headheight}{14pt}

\setCJKmainfont{Microsoft JhengHei}
\setmainfont{Times New Roman}

\definecolor{accent}{RGB}{30, 80, 160}
\definecolor{zhfill}{RGB}{237, 246, 255}
\definecolor{enfill}{RGB}{238, 248, 240}
\definecolor{engreen}{RGB}{40, 120, 70}
\definecolor{timefill}{RGB}{232, 232, 232}

\tcbuselibrary{breakable}

\pagestyle{fancy}
\fancyhf{}
\renewcommand{\headrulewidth}{0.4pt}
\fancyhead[L]{\small\color{gray} 口語報告提示稿（中英對照）}
\fancyhead[R]{\small\color{gray} 黃崇晉 | L16141149 | 數論（一）2026}
\fancyfoot[C]{\small\thepage}

\titleformat{\section}[block]
  {\normalsize\bfseries\color{accent}}{}
  {0pt}{}
  [\vspace{1pt}\titlerule\vspace{3pt}]

\setstretch{1.18}
\setlength{\parskip}{2.0pt}
\setlength{\parindent}{0pt}

\newcommand{\timetag}[1]{\hfill\colorbox{timefill}{\small\color{gray}\textbf{#1}}}

% 投影片標題列
\newcommand{\slide}[2]{%
  \vspace{2pt}%
  \noindent{\small\bfseries\color{accent} ◆ 第 #1 頁\;|\;#2}\par\vspace{1pt}%
}

% 中文口語框
\newenvironment{zh}{%
  \begin{tcolorbox}[colback=zhfill, colframe=accent!45, breakable,
    left=5pt, right=5pt, top=2.5pt, bottom=2.5pt, boxrule=0.4pt,
    title={\footnotesize\bfseries 中文口語}, fonttitle=\color{white},
    coltitle=white, colbacktitle=accent]\small
}{\end{tcolorbox}\vspace{-1pt}}

% 英文口語框
\newenvironment{en}{%
  \begin{tcolorbox}[colback=enfill, colframe=engreen!45, breakable,
    left=5pt, right=5pt, top=2.5pt, bottom=2.5pt, boxrule=0.4pt,
    title={\footnotesize\bfseries English Script}, fonttitle=\color{white},
    coltitle=white, colbacktitle=engreen]\small\itshape
}{\end{tcolorbox}\vspace{1pt}}

\begin{document}

% ── Title ─────────────────────────────────────────────────────────────
\begin{center}
  {\LARGE\bfseries\color{accent} LWE 與 Ring-LWE — 口語報告提示稿}\\[2pt]
  {\large\bfseries Bilingual Speaker Script（中英對照，對應 presentation.pdf 20 頁）}\\[1pt]
  {\small\itshape\color{accent} 理想格密碼學的數論結構 · 約 40 分鐘}\\[5pt]
  \rule{\linewidth}{0.8pt}\\[2pt]
  {\small 黃崇晉（Chung-Chin Huang）\;|\; L16141149 \;|\; 數論（一）\;|\; 2026 春季，國立成功大學}
\end{center}

{\small\color{gray}\itshape 使用說明：左欄頁碼對應投影片；中文框與英文框內容平行，可擇一語言通講，或中英穿插。框內方括號 [\,] 為語氣／停頓提示。}

%% ══════════════════════════════════════════════════════════════════════
\section{開場與大綱 \timetag{約 1.5 分鐘}}

\slide{1}{封面 — LWE 與 Ring-LWE}
\begin{zh}
各位老師、同學好。本報告題目為「LWE 與 Ring-LWE：理想格密碼學的數論結構」。
整場報告的主線是一條轉換路徑：從\textbf{幾何數論}出發，經過\textbf{分圓域與理想格}，
進入今天的重點\textbf{LWE 與 Ring-LWE}，最後談\textbf{後量子標準}。
核心問題只有一個——當量子電腦能破解 RSA 與橢圓曲線後，密碼學該建立在什麼數學困難之上？答案是「格」。
\end{zh}
\begin{en}
Good [morning / afternoon], everyone. The title of this talk is ``LWE and Ring-LWE: the number-theoretic structure of ideal-lattice cryptography.''
The storyline is a single chain: starting from the \emph{geometry of numbers}, passing through \emph{cyclotomic fields and ideal lattices}, into today's focus, \emph{LWE and Ring-LWE}, and ending with the \emph{post-quantum standards}.
There is one central question: once quantum computers break RSA and elliptic curves, on what mathematical hardness should cryptography stand? The answer is: lattices.
\end{en}

\slide{2}{報告大綱與時間配置}
\begin{zh}
報告分四部分。Part~0 用約十分鐘快速鋪陳六個預備概念，每個概念一頁。
Part~1 與 Part~2 是重點，分別講 LWE 與 Ring-LWE，約佔二十五分鐘。
最後 Part~3 用五分鐘談應用與開放問題。重點會放在「為什麼需要 Ring-LWE」以及「環結構如何同時帶來效率與爭議」。
\end{zh}
\begin{en}
The talk has four parts. Part~0 spends about ten minutes laying out six preliminary concepts, one slide each.
Parts~1 and~2 are the core—LWE and then Ring-LWE—taking about twenty-five minutes.
Finally, Part~3 spends five minutes on applications and open problems. The emphasis is on \emph{why} Ring-LWE is needed, and how its ring structure brings both efficiency and controversy.
\end{en}

%% ══════════════════════════════════════════════════════════════════════
\section{Part 0 — 預備概念 \timetag{約 10 分鐘 · 六頁各約 1.5 分鐘}}

\slide{3}{格（Lattice）}
\begin{zh}
首先是格的定義。取 $\mathbb{R}^n$ 中 $n$ 個線性獨立向量，它們的\textbf{整數}線性組合就構成一個秩 $n$ 的格。
請注意關鍵字是「整數」係數，這讓格成為離散的點陣。基底矩陣行列式的絕對值稱為覆積，等於基本域的體積，且與選哪組基底無關。
最重要的直覺是：同一個格有無窮多組基底，有「好基底」（又短又接近正交）也有「壞基底」——這個落差正是之後所有困難問題的來源。
\end{zh}
\begin{en}
We begin with the definition of a lattice. Take $n$ linearly independent vectors in $\mathbb{R}^n$; their \emph{integer} linear combinations form a rank-$n$ lattice.
The keyword is ``integer'' coefficients, which makes the lattice a discrete grid of points. The absolute value of the basis determinant is the covolume, equal to the volume of the fundamental domain, and it is independent of the chosen basis.
The key intuition: one lattice has infinitely many bases—``good'' ones (short, nearly orthogonal) and ``bad'' ones. That gap is exactly the source of every hard problem to come.
\end{en}

\slide{4}{Minkowski 第一定理（1896）}
\begin{zh}
接下來是 Minkowski 第一定理。對一個中心對稱凸體，只要它的體積超過 $2^n$ 乘上覆積，就必定包含一個非零格點。
證明的構想是對二分之一的 $S$ 套用 Blichfeldt 鴿巢論證，逼出兩個平移拷貝重疊，再用對稱與凸性把差向量拉回 $S$。
這裡最該強調的是「非構造性」三個字：定理只保證短向量\textbf{存在}，卻完全不告訴我們怎麼找到它。這個「存在但難尋」的落差，就是 SVP 困難性的幾何根源。
\end{zh}
\begin{en}
Next is Minkowski's first theorem. For a centrally symmetric convex body, as soon as its volume exceeds $2^n$ times the covolume, it must contain a nonzero lattice point.
The proof idea: apply Blichfeldt's pigeonhole argument to one-half of $S$, force two translated copies to overlap, then use symmetry and convexity to pull the difference vector back into $S$.
What to stress is the word \emph{non-constructive}: the theorem guarantees a short vector \emph{exists}, but says nothing about how to find it. That gap—``exists but hard to locate''—is the geometric root of the hardness of SVP.
\end{en}

\slide{5}{分圓域（Cyclotomic Field）}
\begin{zh}
現在把場景換到數論。本報告固定取 $n$ 為 $2$ 的冪，這時第 $2n$ 個分圓多項式恰好是 $x^n+1$，在有理數上不可約。
對應的整數環是 Dedekind 整環，保證每個理想都有唯一的質理想分解——這是後面所有結構的基礎。
透過標準嵌入，把數域中的元素送進 $\mathbb{C}^n$，而且環的乘法剛好對應到逐分量乘法。
選 $x^n+1$ 的三個理由：不可約性最乾淨、$x^n$ 同餘於 $-1$ 帶來負循環卷積、最有利於後面的 NTT。
\end{zh}
\begin{en}
Now we move to number theory. Throughout, we take $n$ to be a power of two; then the $2n$-th cyclotomic polynomial is exactly $x^n+1$, irreducible over the rationals.
Its ring of integers is a Dedekind domain, which guarantees unique factorization of every ideal into prime ideals—the foundation for everything that follows.
Through the canonical embedding, elements of the field map into $\mathbb{C}^n$, and ring multiplication corresponds to component-wise multiplication.
Three reasons for choosing $x^n+1$: cleanest irreducibility, the relation $x^n \equiv -1$ giving negacyclic convolution, and maximal friendliness to the NTT later on.
\end{en}

\slide{6}{理想格結構（Ideal Lattice）}
\begin{zh}
這一頁是整場報告的橋樑。透過標準嵌入，數域中的每個理想都對應到一個格，稱為理想格。
最關鍵的公式是：理想格的覆積等於理想範數乘上根號判別式。換句話說，\textbf{算術上的「理想範數」被精確翻譯成幾何上的「格體積」}。
理想格其實只是一般格的特例，但多了環的乘法結構。這個額外結構是一把雙面刃：它是 Ring-LWE 效率提升的根源，同時也是安全爭議的來源——因為理想格可能比一般格更容易攻擊。
\end{zh}
\begin{en}
This slide is the bridge of the whole talk. Via the canonical embedding, every ideal in the number field maps to a lattice, called an ideal lattice.
The key formula: the covolume of an ideal lattice equals the ideal norm times the square root of the discriminant. In other words, the arithmetic ``ideal norm'' is translated precisely into the geometric ``lattice volume.''
An ideal lattice is just a special case of a general lattice, but with extra ring-multiplication structure. That extra structure is double-edged: it is the source of Ring-LWE's efficiency, and also the source of security concerns—because ideal lattices might be easier to attack than arbitrary ones.
\end{en}

\slide{7}{最短向量問題 SVP}
\begin{zh}
有了格，就能定義第一個核心難題：最短向量問題。給定一組\textbf{壞}基底，要找出最短的非零格點，或它的長度。
這個問題是 NP-hard，目前最佳演算法是指數時間 $2^{\Theta(n)}$；即使用量子電腦，最佳也只是 $2^{0.265n}$，仍是指數級。
這就是為什麼格密碼被認為能抵抗量子攻擊。記住一句話：LWE 與 Ring-LWE 的安全性，最終都歸約到最壞情況下的 SVP。
\end{zh}
\begin{en}
With lattices in place, we can define the first core hard problem: the Shortest Vector Problem. Given a \emph{bad} basis, find the shortest nonzero lattice vector, or its length.
This problem is NP-hard; the best algorithms run in exponential time $2^{\Theta(n)}$. Even on a quantum computer, the best is $2^{0.265n}$—still exponential.
That is precisely why lattice cryptography is believed to resist quantum attacks. Keep one sentence in mind: the security of LWE and Ring-LWE ultimately reduces to worst-case SVP.
\end{en}

\slide{8}{最近向量問題 CVP}
\begin{zh}
第二個難題是最近向量問題：給定一個不一定在格上的目標點，要找離它最近的格點。CVP 一般比 SVP 更難。
其中有個特例叫 BDD，有界距離解碼——已知目標點離某個格點\textbf{非常近}，要把那個格點還原出來，而且解是唯一的。
請特別記住這一頁的最後一句：\textbf{LWE 的本質就是 BDD}。等一下會看到，LWE 樣本裡的 $As$ 是格點、誤差 $e$ 是微小偏移，還原秘密 $s$ 就等於在格上解 BDD。
\end{zh}
\begin{en}
The second hard problem is the Closest Vector Problem: given a target point not necessarily on the lattice, find the nearest lattice point. CVP is generally harder than SVP.
A special case is BDD, Bounded Distance Decoding: the target is known to be \emph{very close} to some lattice point, and we must recover that point—with a unique solution.
Remember the last line of this slide: \emph{LWE is essentially BDD}. We will soon see that in an LWE sample, $As$ is the lattice point, the error $e$ is a tiny offset, and recovering the secret $s$ is exactly solving BDD on a lattice.
\end{en}

%% ══════════════════════════════════════════════════════════════════════
\section{Part 1 — LWE（重點）\timetag{約 12 分鐘}}

\slide{9}{LWE：問題設定（Regev 2005）}
\begin{zh}
進入重點。LWE，Learning With Errors，由 Regev 在 2005 年提出。設定是：固定維度 $n$、模數 $q$、一個窄的高斯誤差分布。
秘密是一個向量 $s$。每筆樣本隨機取一個向量 $a$，計算 $a$ 與 $s$ 的內積，再\textbf{加上一個微小誤差} $e$，模 $q$ 後輸出。
關鍵直覺在這裡：如果\textbf{沒有}誤差，只要收集 $n$ 條方程，用高斯消去法就能解出 $s$。但一旦加入微小誤差，線性代數整個失效——這就是困難的來源。誤差必須恰到好處：小到能正確解密、大到足以保證安全。
\end{zh}
\begin{en}
Now the core. LWE—Learning With Errors—was introduced by Regev in 2005. The setup fixes a dimension $n$, a modulus $q$, and a narrow Gaussian error distribution.
The secret is a vector $s$. Each sample draws a random vector $a$, computes the inner product of $a$ and $s$, then \emph{adds a tiny error} $e$, and outputs the result modulo $q$.
Here is the key intuition: \emph{without} the error, collecting $n$ equations and running Gaussian elimination recovers $s$ immediately. But once a tiny error is added, linear algebra breaks down entirely—that is the source of hardness. The error must be just right: small enough for correct decryption, large enough for security.
\end{en}

\slide{10}{LWE：判定版本與搜尋版本}
\begin{zh}
LWE 有兩個版本。搜尋版：給很多筆樣本，把秘密 $s$ 解出來。判定版：要區分「真正的 LWE 樣本」和「完全均勻的隨機串」。
這兩個版本在質數 $q$ 時是多項式時間等價的，Regev 給出了 search-to-decision 的歸約。
為什麼加密用判定版？因為判定版的困難，直接保證\textbf{密文看起來和隨機串無法區分}，這就是語意安全 IND-CPA。
舉 Regev 加密為例：公鑰就是 $A$ 和 $b$，加密一個位元時取若干筆樣本求和，並把訊息藏在 $q$ 的一半這個位置；解密時看內積靠近 $0$ 還是 $q/2$。
\end{zh}
\begin{en}
LWE comes in two versions. Search: given many samples, recover the secret $s$. Decision: distinguish ``genuine LWE samples'' from ``uniformly random strings.''
For prime $q$ the two are polynomial-time equivalent; Regev gave the search-to-decision reduction.
Why does encryption use the decision version? Because its hardness directly guarantees that \emph{ciphertexts are indistinguishable from random}—that is semantic security, IND-CPA.
Take Regev encryption as an example: the public key is $A$ and $b$; to encrypt a bit, sum a subset of samples and hide the message near $q/2$; to decrypt, check whether the inner product is close to $0$ or to $q/2$.
\end{en}

\slide{11}{LWE：與格的連結與安全歸約}
\begin{zh}
現在把 LWE 翻譯成格的語言。考慮 $q$-ary 格，也就是所有形如 $As$ 模 $q$ 的點所成的格。
那麼 $b = As + e$ 就是一個\textbf{靠近格點 $As$ 的目標點}，誤差 $e$ 正好是那個小偏移。所以還原 $s$，就等於在這個格上解 BDD，呼應前面 CVP 那一頁。
更重要的是 Regev 的量子歸約：他證明了「破解平均情況的 LWE」可以推出「解最壞情況的 GapSVP 或 SIVP」。
這句話的份量在於——安全性不是建立在某些隨機實例上，而是建立在\textbf{最難的格實例}上。Peikert 在 2009 年補上了針對 GapSVP、大 $q$ 的古典歸約。
\end{zh}
\begin{en}
Now translate LWE into lattice language. Consider the $q$-ary lattice—the lattice of all points of the form $As$ modulo $q$.
Then $b = As + e$ is a \emph{target close to the lattice point $As$}, with the error $e$ as the small offset. So recovering $s$ is solving BDD on this lattice—echoing the CVP slide.
More importantly, Regev's quantum reduction proves that ``breaking average-case LWE'' implies ``solving worst-case GapSVP or SIVP.''
The weight of this statement: security rests not on some random instances, but on the \emph{hardest} lattice instances. Peikert in 2009 added a classical reduction for GapSVP with large $q$.
\end{en}

\slide{12}{LWE：困難性與效率瓶頸}
\begin{zh}
LWE 的優點是安全歸約乾淨、不需要任何額外的代數結構假設，所以它是最保守、最讓人放心的格密碼基礎。
但它有一個致命的效率瓶頸：公鑰是一個 $m$ 乘 $n$ 的矩陣，大小是 $n$ 平方等級；每次運算是一般的矩陣向量乘，也是 $n$ 平方。實用安全要求 $n$ 大約上千，金鑰動輒好幾 MB，難以部署。
於是核心問題浮現：能不能在\textbf{不犧牲安全歸約}的前提下，把 $n$ 平方降到接近線性？答案是給秘密和樣本「環結構」——一個環元素就能一次提供 $n$ 條方程。這就引出 Ring-LWE。
\end{zh}
\begin{en}
LWE's advantage is a clean security reduction with no extra algebraic assumptions, making it the most conservative and reassuring foundation for lattice cryptography.
But it has a fatal efficiency bottleneck: the public key is an $m$-by-$n$ matrix, of size order $n$ squared; each operation is a generic matrix-vector product, also $n$ squared. Practical security needs $n$ around a thousand, so keys reach several megabytes—hard to deploy.
Hence the central question: can we cut $n$ squared down to nearly linear \emph{without} sacrificing the security reduction? The answer is to give the secret and samples a ``ring structure''—a single ring element provides $n$ equations at once. This leads us to Ring-LWE.
\end{en}

%% ══════════════════════════════════════════════════════════════════════
\section{Part 2 — Ring-LWE（重點）\timetag{約 13 分鐘}}

\slide{13}{Ring-LWE：分圓環設定（LPR 2010）}
\begin{zh}
Ring-LWE 由 Lyubashevsky、Peikert、Regev 在 2010 年提出。把 LWE 的整個場景搬進分圓環裡。
對照來看：原本的向量空間 $\mathbb{Z}_q^n$ 換成環 $R_q$；原本的內積換成\textbf{環乘法}，也就是負循環卷積。
樣本變成：隨機取一個環元素 $a$，輸出 $a$ 乘 $s$ 再加誤差 $e$。妙處在於：單一個 $a$ 有 $n$ 個係數，它一次就攜帶了 $n$ 條方程，所以金鑰從 $n$ 平方掉到 $n$。
代價是安全性現在歸約到\textbf{理想格}上的 SVP，而不是任意格。接下來三頁說明環結構如何讓運算加速、又如何讓歸約成立。
\end{zh}
\begin{en}
Ring-LWE was introduced by Lyubashevsky, Peikert, and Regev in 2010. It moves the entire LWE setting into a cyclotomic ring.
By analogy: the vector space $\mathbb{Z}_q^n$ becomes the ring $R_q$; the inner product becomes \emph{ring multiplication}, i.e.\ negacyclic convolution.
A sample becomes: draw a random ring element $a$, output $a$ times $s$ plus error $e$. The beauty: a single $a$ has $n$ coefficients, so it carries $n$ equations at once, and the key drops from $n$ squared to $n$.
The price is that security now reduces to SVP on \emph{ideal} lattices, not arbitrary ones. The next three slides show how the ring structure speeds up computation and makes the reduction work.
\end{en}

\slide{14}{Ring-LWE：質理想分解與 $R_q \cong \mathbb{F}_q^n$}
\begin{zh}
要讓乘法變快，先看 $q$ 在環裡怎麼分解。對不整除 $2n$ 的質數 $q$，$q$ 的分解完全由 $x^n+1$ 模 $q$ 的因式分解決定。
最理想的情況是\textbf{完全分裂}，它發生的充要條件非常漂亮：$q$ 同餘於 $1$ 模 $2n$。
這時 $x^n+1$ 在模 $q$ 下裂成 $n$ 個相異的一次因式，由中國剩餘定理，整個環 $R_q$ 同構於 $n$ 份 $\mathbb{F}_q$ 的直積，也就是 $\mathbb{F}_q^n$。
這個同構把「多項式乘法」變成「逐分量乘法」——這正是下一頁 NTT 的代數根源。所以實作上參數永遠選 $q$ 同餘 $1$ 模 $2n$。
\end{zh}
\begin{en}
To make multiplication fast, first see how $q$ factors in the ring. For a prime $q$ not dividing $2n$, the factorization of $q$ is fully determined by the factorization of $x^n+1$ modulo $q$.
The ideal case is \emph{complete splitting}, and its necessary-and-sufficient condition is beautifully simple: $q$ congruent to $1$ modulo $2n$.
Then $x^n+1$ splits modulo $q$ into $n$ distinct linear factors, and by the Chinese Remainder Theorem the whole ring $R_q$ is isomorphic to a product of $n$ copies of $\mathbb{F}_q$, that is $\mathbb{F}_q^n$.
This isomorphism turns ``polynomial multiplication'' into ``component-wise multiplication''—the algebraic root of the NTT on the next slide. So in practice the parameter is always chosen with $q$ congruent to $1$ modulo $2n$.
\end{en}

\slide{15}{Ring-LWE：NTT — $O(n\log n)$ 的環乘法}
\begin{zh}
數論變換 NTT，就是上一頁那個求值同構的演算法化。把多項式 $a$ 送到它在 $n$ 個單位根上的取值，乘法就變成逐分量相乘，乘完再用反變換插值回來。
直接算多項式乘法要 $n$ 平方；但用 Cooley-Tukey 蝴蝶運算，依奇偶指標遞迴拆解，配合 $\omega^n = -1$ 這個負循環性質，遞迴式解出來是 $n \log n$。
這就是 Ring-LWE 效率的引擎。也正因為 $x^n+1$ 給出 $x^n \equiv -1$，乘法是\textbf{負循環卷積}——超過次數 $n$ 的項繞回來時要變號。選 $2$ 的冪分圓環，就是為了讓這個運算最乾淨。
\end{zh}
\begin{en}
The Number-Theoretic Transform, NTT, is the algorithmic form of that evaluation isomorphism. Send a polynomial $a$ to its values at the $n$ roots of unity; multiplication becomes component-wise, then an inverse transform interpolates back.
Direct polynomial multiplication costs $n$ squared; but the Cooley--Tukey butterfly recursively splits by parity of indices, and with the negacyclic property $\omega^n = -1$, the recurrence solves to $n \log n$.
This is the engine of Ring-LWE's efficiency. And precisely because $x^n+1$ gives $x^n \equiv -1$, multiplication is \emph{negacyclic convolution}—terms wrapping past degree $n$ flip sign. Choosing a power-of-two cyclotomic ring keeps this operation cleanest.
\end{en}

\slide{16}{Ring-LWE：對偶格與誤差的自然定義域}
\begin{zh}
這一頁比較技術，但對安全歸約的嚴密性很關鍵。先算差異理想：它由分圓多項式的導數生成，在這裡剛好等於主理想 $(n)$。
在跡配對的意義下，整數環的對偶格就是\textbf{逆差理想}，這裡恰好等於 $n$ 分之一倍的整數環。驗證的關鍵是 $\zeta$ 的次方的跡：只有零次方的跡是 $n$，其餘都是零。
密碼學上的意義是：嚴格的 Ring-LWE，誤差不是定義在 $R_q$ 上，而是定義在對偶格上。唯有用對偶格，誤差分布在跡配對下才有\textbf{代數不變性}，安全歸約才會嚴密。在 $2$ 的冪分圓環裡，這個修正只是個 $n$ 分之一的縮放，實作很簡單。
\end{zh}
\begin{en}
This slide is more technical but crucial for the rigor of the reduction. First compute the different ideal: generated by the derivative of the cyclotomic polynomial, here equal to the principal ideal $(n)$.
Under the trace pairing, the dual lattice of the ring of integers is the \emph{codifferent}, which here equals one-over-$n$ times the ring of integers. The key to verifying is the trace of powers of $\zeta$: only the zeroth power has trace $n$, all others are zero.
The cryptographic meaning: in rigorous Ring-LWE, the error lives not on $R_q$ but on the dual lattice. Only with the dual lattice does the error distribution gain \emph{algebraic invariance} under the trace pairing, making the reduction tight. In a power-of-two cyclotomic ring this correction is just a one-over-$n$ scaling—easy to implement.
\end{en}

\slide{17}{Ring-LWE：歸約鏈與 LWE 對照（小結）}
\begin{zh}
Part~2 小結。安全歸約鏈是：最壞情況的 ideal-SVP，量子多項式時間歸約到平均情況的 Ring-LWE，再推出破解 Kyber、Dilithium 的困難。
看這張對照表：LWE 金鑰 $n$ 平方、乘法 $n$ 平方、建立在任意格上；Ring-LWE 金鑰只要 $n$、乘法 $n \log n$、建立在理想格上。
一句話總結：分圓環、CRT、NTT、對偶格這一整套代數結構，讓 Ring-LWE 在維持與 LWE 相近的安全保證下，把成本從 $n$ 平方壓到 $n \log n$。代價就是多了一個「理想格」的假設——這正是最後開放問題的核心。
\end{zh}
\begin{en}
Summary of Part~2. The reduction chain is: worst-case ideal-SVP reduces, in quantum polynomial time, to average-case Ring-LWE, which implies the hardness of breaking Kyber and Dilithium.
Look at the comparison table: LWE has key size $n$ squared, multiplication $n$ squared, based on arbitrary lattices; Ring-LWE has key size only $n$, multiplication $n \log n$, based on ideal lattices.
In one sentence: the whole algebraic toolkit—cyclotomic ring, CRT, NTT, dual lattice—lets Ring-LWE keep security comparable to LWE while cutting cost from $n$ squared to $n \log n$. The price is one extra ``ideal-lattice'' assumption—which is exactly the heart of the final open problems.
\end{en}

%% ══════════════════════════════════════════════════════════════════════
\section{Part 3 — 應用與開放問題 \timetag{約 5 分鐘}}

\slide{18}{應用：NIST 後量子標準（2024.08）}
\begin{zh}
談應用。2024 年 8 月，NIST 正式發布後量子標準，其中兩個都建立在格上。
FIPS~203，也就是 Kyber，做金鑰封裝，基礎是 Module-LWE——它是 Ring-LWE 在小矩陣上的推廣。Kyber-512 取 $n=256$、$q=3329$，攻擊難度約 $2$ 的 $118$ 次方。
FIPS~204，Dilithium，做數位簽章，基礎是 Module-LWE 加 Module-SIS。
為什麼用 Module 而不純用 Ring？因為 Module-LWE 是個折衷：用一個小矩陣，介於最快但結構最多的 Ring-LWE，和最保守的 LWE 之間，既降低結構風險、又還能用 NTT 加速。整條安全鏈是：破解密碼等價於解理想格上的 BDD，而它又被 ideal-SVP 的困難所保證。
\end{zh}
\begin{en}
On applications. In August 2024, NIST released its post-quantum standards, two of which are lattice-based.
FIPS~203, Kyber, does key encapsulation, based on Module-LWE—a generalization of Ring-LWE over small matrices. Kyber-512 uses $n=256$, $q=3329$, with attack hardness about $2$ to the $118$.
FIPS~204, Dilithium, does digital signatures, based on Module-LWE plus Module-SIS.
Why Module rather than pure Ring? Because Module-LWE is a compromise: using a small matrix, it sits between the fastest, most-structured Ring-LWE and the most conservative LWE—reducing structural risk while still allowing NTT acceleration. The full security chain: breaking the cipher is equivalent to solving BDD on an ideal lattice, which in turn is guaranteed by the hardness of ideal-SVP.
\end{en}

\slide{19}{開放問題}
\begin{zh}
最後是四個開放問題。第一，也是最重要的：理想格上的 SVP，是否真的不比一般格容易？Cramer 等人在 2016 年對某些分圓環給出了量子加速。如果哪天 ideal-SVP 被證明容易，Ring-LWE 的安全保證就會被削弱。
第二，非分圓域的 Ring-LWE：現有歸約依賴分圓域的交換 Galois 結構，換成非交換的情況還不知道行不行。
第三，古典歸約：Regev 的歸約是量子的，純古典的最壞到平均歸約至今沒找到。
第四，量子下界：目前最好的量子演算法是 $2$ 的 $0.265n$ 次方，但無條件的指數下界一直沒被證明。這些都是目前活躍的研究前沿。
\end{zh}
\begin{en}
Finally, four open problems. First, and most important: is SVP on ideal lattices truly no easier than on general lattices? Cramer et al.\ in 2016 gave a quantum speedup for certain cyclotomic rings. If ideal-SVP is ever shown to be easy, Ring-LWE's security guarantee weakens.
Second, Ring-LWE over non-cyclotomic fields: existing reductions rely on the abelian Galois structure of cyclotomic fields; whether they extend to non-abelian cases is unknown.
Third, classical reductions: Regev's reduction is quantum, and a purely classical worst-case-to-average-case reduction has not been found.
Fourth, quantum lower bounds: the best quantum algorithm is $2$ to the $0.265n$, yet an unconditional exponential lower bound remains unproven. These are all active research frontiers.
\end{en}

\slide{20}{結尾與致謝}
\begin{zh}
總結整場報告：從格與 Minkowski 定理的幾何直覺出發，看到分圓域如何把理想變成理想格；再進入 LWE 與它的 BDD 本質，最後看 Ring-LWE 如何用環結構把效率推到實用等級，並成為今天 NIST 標準的數學基礎。
核心訊息只有一句：\textbf{格上「存在但難尋」的短向量，撐起了整個後量子密碼學}。感謝聆聽，歡迎提問。
\end{zh}
\begin{en}
To wrap up: starting from the geometric intuition of lattices and Minkowski's theorem, we saw how cyclotomic fields turn ideals into ideal lattices; then we entered LWE and its BDD essence; and finally how Ring-LWE uses ring structure to push efficiency to practical levels, becoming the mathematical foundation of today's NIST standards.
The core message is one sentence: \emph{the ``exists-but-hard-to-find'' short vector of a lattice supports the whole of post-quantum cryptography}. Thank you for listening; I welcome your questions.
\end{en}

\vspace{4pt}
\noindent\rule{\linewidth}{0.4pt}
\begin{center}\small\color{gray}
  對應投影片：presentation.pdf（20 頁）\quad|\quad 時間：Part 0 \textasciitilde10 分 + Part 1\,/\,2 \textasciitilde25 分 + Part 3 \textasciitilde5 分 $\approx$ \textbf{40 分鐘}\\[2pt]
  O.\ Regev, \textit{J.\ ACM} \textbf{56}(6), 2009.\;
  V.\ Lyubashevsky--C.\ Peikert--O.\ Regev, \textit{J.\ ACM} \textbf{60}(6), 2013.\;
  H.\ Minkowski, 1896.\; J.\ Neukirch, Springer, 1999.\; NIST FIPS 203/204, 2024.
\end{center}

\end{document}
